% =====================================================================
% FAIMR: Journal Extension — Applied Intelligence (Springer)
% Document class: sn-jnl (Springer Nature LaTeX template, smallcondensed)
% Required files: sn-article.cls, sn-basic.bst
% Download from: https://www.springernature.com/gp/authors/campaigns/latex-author-support
% =====================================================================

\documentclass[smallcondensed]{sn-jnl}

\usepackage{amsmath,amssymb,amsfonts}
\usepackage{amsthm}
\usepackage{algorithmic}
\usepackage{graphicx}
\usepackage{float}
\usepackage{booktabs}
\usepackage{multirow}
\usepackage{url}
\def\UrlFont{\rmfamily}

% hyperref is loaded by sn-jnl — configure via \hypersetup only
\hypersetup{colorlinks=true,linkcolor=black,citecolor=black,urlcolor=black}

% Theorem environments
\newtheorem{theorem}{Theorem}
\newtheorem{proposition}{Proposition}
\newtheorem{definition}{Definition}
\newtheorem{empiricalobs}{Empirical Observation}

\begin{document}

\title[FAIMR: Fairness-Aware Resume Ranking]{FAIMR: A Fairness-Aware and Explainable Multi-Signal Framework for Automated Resume Ranking}

\author*[1]{\fnm{Arnav} \sur{Gupta}}\email{arnav.g1010@gmail.com}
\author[1]{\fnm{Virendrakumar A.} \sur{Dhotre}}\email{virendradhotre@gmail.com}

\affil*[1]{\orgdiv{Department of Computer Science and Engineering (AI \& ML)},
  \orgname{Vishwakarma Institute of Technology Pune},
  \orgaddress{\city{Pune}, \postcode{411037}, \state{Maharashtra}, \country{India}}}

\maketitle

\begin{abstract}
Automated Resume Screening (ARS) systems that rely on a single similarity signal suffer from keyword dilution, demographic bias, and decision opacity. We propose FAIMR (Fairness-Aware Explainable Multi-Signal Resume Re-ranking), a framework that fuses Sentence-BERT contextual embeddings, TF-IDF lexical similarity, and skill-graph coverage into a gradient-boosted Learning-to-Rank ensemble. We prove that the three signal families occupy complementary representational subspaces and that their non-linear fusion captures feature interactions inaccessible to any single representation. A post-hoc Fairness-Constrained Re-ranking (FCR) algorithm enforces the four-fifths rule of Adverse Impact while provably minimizing rank displacement, and a counterfactual explainer identifies the minimal skill acquisition path to invert a rejection via greedy submodular optimization with formal approximation guarantees. Evaluation on 26,020 domain-stratified resume--job pairs with verified feature--label independence ($r_{\mathrm{pb}}{=}0.14$) shows FAIMR achieves AUC $0.719 \pm 0.010$, outperforming BM25 ($0.604$), Doc2Vec ($0.596$), BERT Cross-Encoder ($0.585$), MPNet-SBERT ($0.656$), and Jaccard Skill Matching ($0.581$), with each increment statistically significant ($p < 0.002$). The FCR module achieves Adverse Impact Ratio $0.979$ on unbiased data and restores AIR~$\geq 0.8$ at all synthetic skew levels. The counterfactual generator explains 5,158 rejected candidates with non-empty skill deficiencies, averaging $1.10$ skills with median latency $0.62$\,ms and $100\%$ greedy-optimal match rate. Feature importance analysis confirms semantic similarity ($\approx 35\%$) and skill coverage ($\approx 28\%$) as dominant signals, validating the theoretical complementarity argument.
\end{abstract}

\keywords{Learning-to-Rank, Algorithmic Fairness, Natural Language Processing, Resume Screening, Explainable AI, Counterfactual Reasoning}

% =====================================================================
\section{Introduction}
\label{sec:intro}

The digital transformation of human resources has produced an unprecedented asymmetry between the volume of job applications and the capacity of hiring teams to evaluate them. Organizations routinely receive thousands of applications per position, making manual screening economically infeasible~\cite{chen2023automated,kessler2018algorithm}. This has driven widespread adoption of Automated Resume Screening (ARS) systems that rank candidates by relevance to a given Job Description~(JD).

Early ARS systems relied on lexical techniques such as Term Frequency--Inverse Document Frequency (TF-IDF) and BM25~\cite{robertson2009probabilistic}. While computationally efficient, these methods suffer from the vocabulary mismatch problem: semantically equivalent phrases (\textit{e.g.}, ``front-end engineering'' vs.\ ``UI development'') produce disjoint feature vectors, yielding high false-negative rates~\cite{singh2020automated}. The introduction of dense sentence embeddings via Sentence-BERT~\cite{reimers2019sbert} addressed vocabulary mismatch but introduced three new failure modes that constitute the central motivation for this work.

\subsection{Motivating Failure Modes}

\textbf{Failure Mode 1: Keyword Dilution.} Dense embeddings project all tokens into a single fixed-dimensional vector through mean pooling. Consider a JD requiring $n$ skills $\{s_1, \ldots, s_n\}$ and a resume possessing $n{-}1$ of them. The embedding cosine similarity will be approximately $(n{-}1)/n$ of the perfect-match similarity, yet the single missing skill may be a dealbreaker (\textit{e.g.}, a regulatory certification). Formally, let $\mathbf{E}(T) = \frac{1}{|T|}\sum_{t \in T} \mathbf{e}(t)$ denote mean-pooled embeddings. The contribution of any single token $t_k$ to the final vector diminishes as $O(1/|T|)$, making dense representations inherently insensitive to individual skill presence or absence.

\textbf{Failure Mode 2: Algorithmic Bias.} Models trained on historical hiring decisions encode societal biases as statistical regularities. Proxy variables correlated with protected attributes---graduation institution, extracurricular terminology, writing style---produce discriminatory rankings. The U.S. Equal Employment Opportunity Commission's (EEOC) four-fifths rule provides a legal threshold: the selection rate for any protected group must be at least 80\% of the rate for the highest-selected group~\cite{dwork2012fairness,buolamwini2018gender}.

\textbf{Failure Mode 3: Decision Opacity.} The non-linear transformations of deep networks sever the causal link between individual qualifications and ranking scores, preventing actionable feedback. Under emerging regulations (EU AI Act Article~14, NYC Local Law~144), high-risk AI systems in employment must provide meaningful explanations~\cite{wachter2017counterfactual}.

\subsection{Novelty and Contributions}

Existing approaches address these failure modes in isolation: fairness-aware ranking algorithms~\cite{zehlike2017fa} do not integrate skill-level explainability, and counterfactual explanation frameworks~\cite{mothilal2020dice} do not incorporate fairness constraints. FAIMR is, to our knowledge, the \textit{first framework to jointly address all three failure modes} through a unified architecture comprising multi-signal fusion, post-hoc fairness enforcement, and counterfactual skill-gap explanation.

The specific novelties are:

\begin{enumerate}
    \item \textbf{Signal-Theoretic Multi-Signal Fusion.} We prove that semantic, lexical, and skill-graph signals occupy complementary representational subspaces (Section~\ref{sec:theory}), and that their non-linear combination via gradient boosting captures conjunctive feature interactions (\textit{e.g.}, ``high semantic match AND low skill coverage implies false positive'') inaccessible to any linear or single-signal model.

    \item \textbf{Leakage-Free Evaluation Methodology.} We identify and eliminate a label-leakage vulnerability in prior skill-coverage-based pair construction, where the labeling criterion (Skill Coverage Ratio) was simultaneously used as a training feature, and propose domain-match labeling with point-biserial verification.

    \item \textbf{Zero-Cost Fairness Guarantee.} The FCR module operates exclusively on the output ranking, preserving the underlying relevance model's predictive performance while providing a provable AIR~$\geq \tau$ guarantee (Section~\ref{sec:fcr_theory}).

    \item \textbf{Optimal Counterfactual Generation.} We formalize skill-gap explanation as a minimum-cardinality set cover problem and prove that the greedy algorithm achieves a $(1 - 1/e)$-approximation under submodularity, with verified exact optimality on bounded instances (Section~\ref{sec:cf_theory}).
\end{enumerate}

% =====================================================================
\section{Related Work}
\label{sec:related}

\subsection{Document Retrieval for Resume Matching}

Classical information retrieval treats resume matching as asymmetric document retrieval. The BM25 probabilistic model~\cite{robertson2009probabilistic} scores documents using term frequency saturation and document length normalization:
\begin{equation}
    \text{BM25}(q, d) = \sum_{t \in q} \text{IDF}(t) \cdot \frac{f(t,d) \cdot (k_1 + 1)}{f(t,d) + k_1 \cdot (1 - b + b \cdot |d|/\text{avgdl})}
\end{equation}
where $f(t,d)$ is term frequency. While effective for exact matches, BM25 assigns zero score to synonymous terms, producing false negatives. Singh~\textit{et al.}~\cite{singh2020automated} document this limitation in first-generation applicant tracking systems, reporting that BM25-based screening rejects up to 30\% of qualified candidates whose resumes use field-standard synonyms not present in the JD vocabulary.

Le and Mikolov~\cite{le2014doc2vec} introduced Doc2Vec (Paragraph Vectors), learning fixed-length document representations via distributed memory (PV-DM) and distributed bag-of-words (PV-DBOW) objectives. While an advance over bag-of-words, Doc2Vec lacks attention mechanisms and produces context-\textit{independent} embeddings. Its performance is particularly limited on short, structured documents such as job descriptions, where the training signal for paragraph-level coherence is weak relative to longer narrative text.

\subsection{Neural Bi-Encoders and Cross-Encoders}

The Transformer architecture~\cite{vaswani2017attention} and BERT~\cite{devlin2018bert} enabled context-dependent embeddings. Reimers and Gurevych~\cite{reimers2019sbert} proposed SBERT, a Siamese bi-encoder producing sentence embeddings for efficient pairwise comparison. The bi-encoder architecture enables sub-linear retrieval through pre-computed embedding caches, making it practical for large-scale candidate pools. Song~\textit{et al.}~\cite{song2020mpnet} introduced MPNet, combining masked and permuted language modeling, achieving state-of-the-art results on sentence similarity benchmarks.

Cross-encoders~\cite{nogueira2019crossencoder} jointly encode document pairs through full cross-attention, enabling richer interaction modeling at the cost of $O(n^2)$ inference complexity. However, cross-encoders pre-trained on web search (MS MARCO) may not transfer effectively to domain-specific recruitment tasks without fine-tuning. Recent work by~\cite{luo2023hire} demonstrates that domain-adaptive pre-training on recruitment corpora improves cross-encoder performance by 8--12 AUC points on resume--JD matching benchmarks, though this fine-tuning requirement represents a significant operational overhead for organizations seeking off-the-shelf solutions.

\textbf{Gap in Prior Work.} All of the above methods produce a \textit{single} relevance signal. While individual signals capture different aspects of document similarity, none can simultaneously detect semantic alignment, exact terminology matches, and explicit skill coverage. FAIMR addresses this gap through principled multi-signal fusion.

\subsection{Learning-to-Rank}

LTR methods provide a framework for combining heterogeneous relevance signals. Burges~\textit{et al.}~\cite{burges2005learning} established listwise and pairwise LTR objectives. Chen and Guestrin~\cite{chen2016xgboost} introduced XGBoost, a scalable gradient-boosted tree ensemble that learns non-linear feature interactions through iterative functional gradient descent. The key theoretical advantage of tree-based LTR over linear combination (weighted averaging) is the ability to learn \textit{conjunctive rules}: ``IF $S_\text{sem} > 0.8$ AND $C_\text{skill} < 0.2$ THEN irrelevant,'' which captures the keyword dilution failure mode directly. Recent extensions explore neural LTR approaches that combine pre-trained language models with learned feature interactions~\cite{chen2023automated}, though these approaches typically sacrifice interpretability and require substantially more training data than gradient-boosted ensembles.

\subsection{Algorithmic Fairness in Ranking}

Dwork~\textit{et al.}~\cite{dwork2012fairness} formalized individual and group fairness. Hardt~\textit{et al.}~\cite{hardt2016equality} contributed Equalized Odds. In ranking, Zehlike~\textit{et al.}~\cite{zehlike2017fa} proposed FA*IR with statistical representation guarantees. Foulds~\textit{et al.}~\cite{foulds2020intersectional} extended fairness to intersectional subgroups. Buolamwini and Gebru~\cite{buolamwini2018gender} demonstrated significant disparities in commercial ML systems. In the specific context of hiring, recent audits of commercial ATS systems have documented demographic disparities ranging from 12\% to 31\% in selection rate differences~\cite{raghavan2020mitigating}, underscoring the practical urgency of post-processing fairness interventions.

Existing fairness interventions are typically \textit{in-processing} (modifying the training objective) or \textit{post-processing} (adjusting outputs). Post-processing is preferred in recruitment because it preserves the integrity and auditability of the relevance model. FAIMR adopts a post-processing approach with formal displacement guarantees.

\subsection{Counterfactual Explanations}

Wachter~\textit{et al.}~\cite{wachter2017counterfactual} introduced counterfactual explanations for affected individuals. Mothilal~\textit{et al.}~\cite{mothilal2020dice} proposed DiCE for gradient-based diverse counterfactual generation. However, these methods assume \textit{continuous} feature spaces. In resume screening, skills are \textit{discrete} (present/absent), transforming counterfactual generation into a combinatorial optimization problem. FAIMR exploits the bounded cardinality of skill requirements to provide exact solutions via greedy submodular optimization.

\textbf{Limitation of scope.} We note that no prior work has experimentally compared automated skill-gap recommendation methods within an ARS ranking context. Zhang~\textit{et al.}~\cite{zhang2022skillspan} address hard and soft skill extraction from job postings, while Decorte~\textit{et al.}~\cite{decorte2022jobbert} develop job-title understanding via skill-grounded representations; both works focus on skill extraction rather than counterfactual explanation within a ranking framework. A direct experimental comparison with DiCE or other counterfactual methods is not conducted in this work, as adapting continuous-space methods to discrete skill sets requires non-trivial modifications that constitute a separate research contribution. We leave this comparison to future work.

\subsection{Resume Screening in Practice}
\label{sec:related_practice}

Commercial Applicant Tracking Systems (ATS) such as Taleo (Oracle), Workday, and Greenhouse process the majority of enterprise hiring workflows. These systems share a common keyword-matching architecture: JD terms are tokenized and candidate resumes are scored by term presence and frequency, typically using proprietary extensions of TF-IDF or BM25. Published audits and industry reports identify three systemic limitations of these platforms. First, they exhibit severe vocabulary sensitivity: resumes using synonymous terminology (\textit{e.g.}, ``machine learning engineer'' vs.\ ``AI/ML specialist'') may receive substantially lower scores despite equivalent qualifications. Second, they provide minimal transparency to candidates, offering no explanation of why a particular resume was ranked below others. Third, compliance with emerging fairness regulations (EU AI Act, NYC Local Law~144) requires audit trail documentation that most current ATS implementations do not natively support~\cite{raghavan2020mitigating}.

Recent transformer-based ARS systems have begun to address vocabulary sensitivity. Chen~\textit{et al.}~\cite{chen2023automated} survey 47 deep learning approaches to resume screening published between 2018 and 2023, finding that bi-encoder architectures achieve the best accuracy--efficiency trade-off but that fewer than 15\% of surveyed systems incorporate any fairness mechanism. Kessler~\textit{et al.}~\cite{kessler2018algorithm} demonstrate that hybrid lexical-semantic pipelines outperform pure neural approaches on specialized recruitment corpora, a finding that motivates FAIMR's fusion design.

Table~\ref{tab:prior_comparison} presents a structured comparison of representative prior ARS methods against FAIMR across five methodology dimensions: signal type, fairness module, explainability mechanism, label-leakage prevention, and evaluation rigor. The comparison reveals that no prior system addresses all five dimensions simultaneously. Fairness-only systems (\textit{e.g.}, FA*IR~\cite{zehlike2017fa}) provide ranking guarantees but ignore skill-level explanation; explanation-only systems (\textit{e.g.}, DiCE~\cite{mothilal2020dice}) generate counterfactuals but assume unbiased rankings; and multi-signal fusion systems (\textit{e.g.}, neural LTR~\cite{chen2023automated}) improve accuracy but lack post-hoc fairness enforcement. FAIMR is the only framework to integrate all five dimensions within a single deployable system.

\begin{table}[t]
\centering
\caption{Comparison of prior ARS methods and FAIMR across five methodology dimensions. \checkmark = fully addressed; $\sim$ = partially addressed; $\times$ = not addressed. LTR = Learning-to-Rank; CF = Counterfactual; LL = Label Leakage.}
\label{tab:prior_comparison}
\begin{tabular}{llccccc}
\toprule
Method & Signal Type & Fairness & Explainability & LL-Free & Eval Rigor \\
\midrule
BM25~\cite{robertson2009probabilistic} & Lexical & $\times$ & $\times$ & $\sim$ & $\sim$ \\
Doc2Vec~\cite{le2014doc2vec} & Semantic (single) & $\times$ & $\times$ & $\sim$ & $\sim$ \\
SBERT bi-encoder~\cite{reimers2019sbert} & Semantic (single) & $\times$ & $\times$ & $\sim$ & $\checkmark$ \\
FA*IR~\cite{zehlike2017fa} & Any (post-proc.) & $\checkmark$ & $\times$ & $\sim$ & $\sim$ \\
DiCE~\cite{mothilal2020dice} & Continuous & $\times$ & $\checkmark$ & $\sim$ & $\sim$ \\
Neural LTR~\cite{chen2023automated} & Multi-signal & $\times$ & $\times$ & $\sim$ & $\sim$ \\
\textbf{FAIMR (Ours)} & \textbf{Multi-signal} & $\checkmark$ & $\checkmark$ & $\checkmark$ & $\checkmark$ \\
\bottomrule
\end{tabular}
\end{table}

% =====================================================================
\section{Theoretical Foundations}
\label{sec:theory}

This section establishes the theoretical basis for FAIMR's multi-signal architecture and provides formal guarantees for the fairness and explainability modules.

\subsection{Signal Complementarity Analysis}

Let $\mathcal{S} = \{S_\text{sem}, S_\text{lex}, S_\text{skill}\}$ denote the three signal families. We argue that these signals are \textit{informationally complementary}: each captures relevance aspects invisible to the others.

\begin{proposition}[Semantic--Lexical Orthogonality]\label{prop:orth}
Let $\mathcal{V}_\text{sem}$ and $\mathcal{V}_\text{lex}$ denote the representational subspaces of SBERT and TF-IDF respectively. There exist document pairs $(d_1, d_2)$ such that $S_\text{sem}(d_1, d_2) \approx 1$ but $S_\text{lex}(d_1, d_2) \approx 0$, and vice versa.
\end{proposition}

\textit{Proof.} Consider two documents $d_1 =$ ``Python developer with machine learning expertise'' and $d_2 =$ ``ML engineer skilled in Python programming.'' These share zero 3-gram TF-IDF tokens (after stop-word removal) yet have high SBERT similarity due to semantic equivalence. Conversely, $d_3 =$ ``Java developer for banking applications'' and $d_4 =$ ``Java developer for healthcare systems'' share identical lexical structure (high TF-IDF similarity) but describe different domains (lower semantic similarity in context). \hfill $\square$

\begin{proposition}[Skill-Signal Independence]\label{prop:skill}
The Skill Coverage Ratio $C_\text{skill}$ provides information that is conditionally independent of both $S_\text{sem}$ and $S_\text{lex}$ given the true relevance label $y$.
\end{proposition}

\textit{Justification.} $C_\text{skill}$ operates on \textit{extracted entities} (discrete skill names) rather than on the continuous token distribution. Two resumes may achieve identical SBERT and TF-IDF scores against a JD while having different skill coverage---\textit{e.g.}, one resume discusses ``machine learning'' extensively (high semantic/lexical overlap) but lacks the specific required skill ``TensorFlow'' (low coverage). Empirically, the measured Pearson correlation between $S_\text{sem}$ and $C_\text{skill}$ across our 26,020 pairs is $r = 0.31$, confirming moderate independence.

\subsection{Why Non-Linear Fusion Outperforms Linear Combination}

A linear scoring function $F_\text{lin} = \alpha S_\text{sem} + \beta S_\text{lex} + \gamma C_\text{skill}$ cannot express \textit{conjunctive conditions}. Consider the keyword dilution scenario: a resume with $S_\text{sem} = 0.92$ (high semantic match) but $C_\text{skill} = 0.15$ (missing critical skills). A linear model may still score this highly. In contrast, a gradient-boosted tree learns the decision boundary:

\begin{equation}
    \text{IF } S_\text{sem} > 0.8 \text{ AND } C_\text{skill} < 0.3 \implies \hat{y} \approx 0
\end{equation}

This conjunctive rule is precisely the \textit{mandatory-skill veto} that single-signal models cannot implement. More formally:

\begin{theorem}[Expressiveness Advantage]\label{thm:express}
Let $\mathcal{F}_\text{lin}$ denote the class of linear scoring functions on $\mathbb{R}^d$ and $\mathcal{F}_\text{tree}$ denote the class of depth-$D$ decision tree ensembles. For $D \geq 2$, there exist classification problems on $\mathbb{R}^d$ where $\inf_{f \in \mathcal{F}_\text{lin}} \text{err}(f) > 0$ while $\inf_{f \in \mathcal{F}_\text{tree}} \text{err}(f) = 0$.
\end{theorem}

\textit{Proof.} The XOR problem on two features is a well-known example. More relevantly, the ``high similarity but low skill coverage'' pattern in resume matching defines a decision boundary that is a \textit{product} of feature conditions, which lies outside the linear hypothesis class but within the tree class for $D \geq 2$. \hfill $\square$

This theoretical result explains the empirical observation that FAIMR (AUC~$= 0.719$) outperforms every single-signal baseline ($\text{AUC} \leq 0.656$): the baselines are restricted to linear decision boundaries in a one-dimensional score space, while FAIMR's XGBoost ensemble operates in the full five-dimensional feature space with non-linear interactions.

\subsection{Fairness-Constrained Re-ranking: Correctness and Optimality}
\label{sec:fcr_theory}

\begin{theorem}[FCR Termination and Correctness]\label{thm:fcr}
Algorithm~1 (FCR) terminates in at most $\min(|L_{>\!k}^{\text{min}}|, |L_{\leq k}^{\text{maj}}|)$ iterations, where $L_{>\!k}^{\text{min}}$ is the set of minority candidates below cutoff $k$ and $L_{\leq k}^{\text{maj}}$ is the set of majority candidates at or above cutoff $k$. Upon termination, either $\mathrm{AIR}(L', k) \geq \tau$ or no further improvement is possible.
\end{theorem}

\textit{Proof.} Each iteration performs exactly one swap, increasing the count of minority candidates in the top-$k$ by one and decreasing the majority count by one. Since the minority count in the top-$k$ is bounded above by $\min(|L^{\text{min}}|, k)$ and increases monotonically, the algorithm must terminate. The AIR is a monotonically non-decreasing function of the minority selection rate when the majority rate is non-increasing, guaranteeing progress toward the threshold. \hfill $\square$

\begin{proposition}[Displacement Minimality]\label{prop:disp}
The greedy swap strategy---swapping the highest-scored rejected minority with the lowest-scored selected majority---minimizes the per-swap score displacement $|\hat{y}(c^+) - \hat{y}(c^-)|$ among all feasible swaps.
\end{proposition}

\textit{Proof.} By construction, $c^+$ has the highest score among rejected minority candidates, and $c^-$ has the lowest score among selected majority candidates, minimizing $|\hat{y}(c^+) - \hat{y}(c^-)|$. \hfill $\square$

\subsection{Counterfactual Optimality Guarantee}
\label{sec:cf_theory}

\begin{definition}[Score Gain Function]
Define $g: 2^\Delta \to \mathbb{R}$ by $g(\delta) = F(\mathbf{x}_{R \cup \delta}) - F(\mathbf{x}_R)$, where $\Delta = \mathcal{K}_J \setminus \mathcal{K}_R$.
\end{definition}

\setcounter{empiricalobs}{3}
\begin{empiricalobs}[Submodularity of Score Gain]\label{prop:submod}
Under the assumption that the XGBoost model's response to skill-coverage features exhibits diminishing marginal returns, the score gain function $g$ is submodular: for $A \subseteq B \subseteq \Delta$ and $s \notin B$,
\begin{equation}
    g(A \cup \{s\}) - g(A) \geq g(B \cup \{s\}) - g(B)
\end{equation}
\end{empiricalobs}

\textit{Justification.} Adding a skill to a resume with few matched skills increases the coverage ratio $C_\text{skill}$ by $1/|\mathcal{K}_J|$. The XGBoost model's learned response to coverage increases is empirically concave (diminishing returns): moving from 0\% to 20\% coverage has a larger impact on the predicted probability than moving from 60\% to 80\%, because the tree splits capturing low-coverage rejection are sharper than those at higher coverage levels. While exact submodularity depends on the learned tree structure, our empirical verification confirms $100\%$ greedy-optimal match rate, validating the submodularity assumption in practice.

Under exact submodularity, the greedy algorithm provides the following guarantee:

\begin{theorem}[Greedy Approximation Bound~\protect\cite{nemhauser1978submodular}]
For submodular $g$ with $g(\emptyset) = 0$, the greedy algorithm selecting $k$ elements achieves:
\begin{equation}
    g(\delta_\text{greedy}) \geq \left(1 - \frac{1}{e}\right) \cdot \max_{|\delta| = k} g(\delta)
\end{equation}
\end{theorem}

Since our greedy search terminates when $F(\mathbf{x}_{R \cup \delta}) \geq \theta$ (a threshold condition rather than cardinality constraint), the algorithm finds the minimum-cardinality solution under submodularity. We verify this by brute-force enumeration for all instances with $|\Delta| \leq 8$, confirming $100\%$ optimality.

% =====================================================================
\section{Methodology}
\label{sec:method}

Let a Job Description $J$ have text $T_J$ and required skills $\mathcal{K}_J$. A Resume $R$ has text $T_R$ and skills $\mathcal{K}_R$. FAIMR learns $F: (J, R) \to \mathbb{R}$ subject to fairness and explainability constraints.

\subsection{Phase 1: Multi-Signal Feature Extraction}

FAIMR extracts five features across three representational dimensions, each targeting a specific failure mode.

\subsubsection{Semantic Similarity (SBERT) --- Addressing Vocabulary Mismatch}
The \texttt{all-MiniLM-L6-v2} Sentence Transformer projects text into $\mathbb{R}^{384}$ via a 6-layer Transformer with mean pooling:
\begin{equation}
    S_{\mathrm{sem}}(J, R) = \frac{\mathbf{E}(T_J) \cdot \mathbf{E}(T_R)}{\|\mathbf{E}(T_J)\| \cdot \|\mathbf{E}(T_R)\|}
\end{equation}
This captures contextual alignment: ``CI/CD pipeline management'' scores highly against ``DevOps orchestration'' despite zero lexical overlap. However, by Proposition~\ref{prop:orth}, this signal alone is blind to exact terminology.

\subsubsection{Lexical Similarity (TF-IDF) --- Preserving Term Specificity}
A TF-IDF vectorizer fitted on the joint corpus (5,000-term vocabulary) computes:
\begin{equation}
    S_{\mathrm{lex}}(J, R) = \frac{\mathbf{v}(T_J) \cdot \mathbf{v}(T_R)}{\|\mathbf{v}(T_J)\| \cdot \|\mathbf{v}(T_R)\|}
\end{equation}
The IDF weighting $\log(N/df_t)$ naturally upweights rare technical terms (\textit{e.g.}, ``Kubernetes'', ``TensorFlow'') that carry high discriminative value. This signal captures exact terminology matches that dense embeddings dilute.

\subsubsection{Skill-Graph Coverage --- Resolving Keyword Dilution}
Explicit skill sets are extracted via a curated ontology. Three deterministic features are derived:
\begin{equation}
    C_{\mathrm{skill}}(J, R) = \frac{|\mathcal{K}_J \cap \mathcal{K}_R|}{|\mathcal{K}_J|}, \quad
    N_{\mathrm{match}} = |\mathcal{K}_J \cap \mathcal{K}_R|
\end{equation}
\begin{equation}
    O_{\mathrm{kw}}(J, R) = \frac{|\mathrm{Tokens}(T_J) \cap \mathrm{Tokens}(T_R)|}{|\mathrm{Tokens}(T_J)|}
\end{equation}

Crucially, $C_\text{skill}$ directly detects the keyword dilution failure mode: a resume with $S_\text{sem} = 0.9$ but $C_\text{skill} = 0.1$ contains relevant discussion of the domain but lacks specific required competencies. By Proposition~\ref{prop:skill}, this information is not redundant with either $S_\text{sem}$ or $S_\text{lex}$.

\subsection{Phase 2: Gradient-Boosted Learning-to-Rank}

The five features are concatenated into $\mathbf{x} = [S_{\mathrm{sem}},\, S_{\mathrm{lex}},\, C_{\mathrm{skill}},\, N_{\mathrm{match}},\, O_{\mathrm{kw}}]^\top$ and input to an XGBoost classifier~\cite{chen2016xgboost}. The ensemble of $K$ additive trees maps $\mathbf{x}$ to a relevance probability:
\begin{equation}
    \hat{y} = \sigma\!\left(\sum_{k=1}^K f_k(\mathbf{x})\right)
\end{equation}
At each boosting iteration $t$, the regularized objective is minimized via second-order Taylor expansion of the logistic loss:
\begin{equation}\label{eq:xgb}
    \mathcal{L}^{(t)} = \sum_{i=1}^n \left[ g_i\, f_t(\mathbf{x}_i) + \tfrac{1}{2}\, h_i\, f_t^2(\mathbf{x}_i) \right] + \Omega(f_t)
\end{equation}
where $g_i = \partial_{\hat{y}} \ell(y_i, \hat{y})$ and $h_i = \partial^2_{\hat{y}} \ell(y_i, \hat{y})$ are the gradient and Hessian, and $\Omega(f_t) = \gamma T + \frac{1}{2}\lambda\|w\|^2$ regularizes tree complexity (leaves $T$, weights $w$).

By Theorem~\ref{thm:express}, this ensemble learns the conjunctive mandatory-skill veto that single-signal models cannot express. The tree structure naturally discovers splits such as ``$C_\text{skill} < 0.25$ AND $S_\text{sem} > 0.7 \implies$ reject,'' directly encoding the intuition that high semantic similarity does not compensate for missing critical requirements.

\subsection{Phase 3: Fairness-Constrained Re-ranking}
\label{sec:fcr_method}

Let $L$ denote the ranked list, $k$ the selection cutoff, and $A_c \in \{0, 1\}$ the demographic group.

\begin{definition}[Adverse Impact Ratio]
\begin{equation}
    \mathrm{AIR}(L, k) = \frac{\min_a P(\text{Sel} \mid A{=}a)}{\max_a P(\text{Sel} \mid A{=}a)}
\end{equation}
where $P(\text{Sel} \mid A{=}a) = |\{c \in L_{1:k} : A_c = a\}| / |\{c \in L : A_c = a\}|$.
\end{definition}

The FCR algorithm enforces $\mathrm{AIR} \geq \tau = 0.8$:

\vspace{3pt}
\noindent\textbf{Algorithm 1:} Fairness-Constrained Re-ranking
\begin{algorithmic}[1]
\REQUIRE Ranked list $L$ sorted by $\hat{y}$, threshold $\tau$, cutoff~$k$
\STATE $L' \leftarrow L$
\WHILE{$\mathrm{AIR}(L', k) < \tau$}
    \STATE $c^+ \leftarrow \arg\max_{i > k}\{\hat{y}(c_i) : A_{c_i} = \text{minority}\}$
    \STATE $c^- \leftarrow \arg\min_{j \leq k}\{\hat{y}(c_j) : A_{c_j} = \text{majority}\}$
    \IF{$c^+$ and $c^-$ exist}
        \STATE Swap positions of $c^+$ and $c^-$ in $L'$
    \ELSE
        \STATE \textbf{break}
    \ENDIF
\ENDWHILE
\RETURN $L'$
\end{algorithmic}
\vspace{3pt}

By Theorem~\ref{thm:fcr}, this algorithm terminates in finite iterations with guaranteed progress toward the AIR threshold. By Proposition~\ref{prop:disp}, the greedy swap strategy minimizes per-swap displacement. The displacement cost is:
\begin{equation}
    \mathcal{D}(L, L') = \frac{1}{n}\sum_{c} |r_L(c) - r_{L'}(c)|
\end{equation}

\subsection{Phase 4: Counterfactual Skill-Gap Explanations}

For a rejected candidate with $\hat{y}_R < \theta$, we seek the minimum skill acquisition set:
\begin{equation}
    \delta^* = \arg\min_{\delta \subseteq \Delta} |\delta| \quad \text{s.t.} \quad F(\mathbf{x}_{R \cup \delta}) \geq \theta
\end{equation}
where $\Delta = \mathcal{K}_J \setminus \mathcal{K}_R$. The greedy algorithm iteratively adds the skill maximizing marginal score improvement:
\begin{equation}
    s_{t+1} = \arg\max_{s \in \Delta \setminus \delta_t} \left[ F(\mathbf{x}_{R \cup \delta_t \cup \{s\}}) - F(\mathbf{x}_{R \cup \delta_t}) \right]
\end{equation}
Only discrete features ($C_\text{skill}$, $N_\text{match}$) are updated per perturbation; semantic and lexical features remain fixed, reflecting the assumption that acquiring a skill changes qualifications but not writing style. By Empirical Observation~\ref{prop:submod}, the greedy solution is optimal or near-optimal under diminishing returns. For $|\Delta| \leq 8$, brute-force enumeration verifies exact optimality.

\subsection{Implementation Details}
\label{sec:impl}

To support reproducibility and deployment planning, we document the full implementation configuration.

\textbf{Sentence-BERT.} All semantic features use the \texttt{all-MiniLM-L6-v2} checkpoint (HuggingFace \texttt{sentence-transformers} library, version 2.2.2). The model produces 384-dimensional embeddings via a 6-layer Transformer with mean-pooling over token representations. Inference is performed on CPU; no GPU is required. Encoding a 500-token resume takes approximately 18\,ms on a standard laptop CPU (Intel Core i7, 2022).

\textbf{TF-IDF Vectorizer.} Implemented via \texttt{sklearn.feature\_extraction.text.TfidfVectorizer} with \texttt{max\_features=5000}, \texttt{sublinear\_tf=True} (log-scaled term frequencies), English stop-word removal, and L2 normalization. The vocabulary is fitted jointly on the training split of each fold to prevent data leakage from the test set.

\textbf{XGBoost Classifier.} Configuration: \texttt{n\_estimators=200}, \texttt{max\_depth=6}, \texttt{learning\_rate=0.1}, \texttt{subsample=0.8}, \texttt{colsample\_bytree=0.8}, \texttt{scale\_pos\_weight=1} (balanced classes), \texttt{eval\_metric=auc}, early stopping on fold validation set with patience 20. The \texttt{scale\_pos\_weight=1} setting reflects the exactly balanced 50/50 positive/negative dataset.

\textbf{Skill Ontology.} A curated list of 847 technical skills across 24 professional domains, developed by consolidating entries from O*NET (v27.2), LinkedIn Skills taxonomy (Q1 2023 snapshot), and manual expert review for the IT, Engineering, Healthcare, Finance, and Marketing domains most heavily represented in the dataset. Skills are matched case-insensitively with a 3-character minimum stem to reduce inflection sensitivity (\textit{e.g.}, ``Python'' matches ``Python 3,'' ``Python scripting'').

\textbf{FCR Cutoff.} The selection cutoff $k$ is set to 20\% of the pool size per fold ($k = 0.2 \times |\text{fold test set}|$), consistent with realistic top-$k$ shortlisting ratios in enterprise hiring workflows.

\textbf{Hardware and Runtime.} All experiments were conducted on a standard consumer laptop (Intel Core i7-1260P, 16\,GB RAM) with no GPU. The full pipeline---feature extraction, XGBoost training across 5 folds, evaluation, FCR stress testing, and counterfactual generation for all 5,158 candidates with non-empty skill deficiencies---completes in under 4 minutes wall-clock time. Feature extraction (SBERT encoding of 26,020 pairs) accounts for approximately 60\% of this time; XGBoost training and evaluation account for the remainder. This runtime profile makes FAIMR practically deployable without specialized infrastructure.

% =====================================================================
\section{Experimental Design}
\label{sec:experiment}

\subsection{Dataset}

\begin{table}[t]
\centering
\caption{Dataset statistics}
\label{tab:dataset}
\begin{tabular}{lr}
\toprule
Parameter & Value \\
\midrule
Unique Job Descriptions & 2,602 \\
Unique Resumes & 4,967 \\
Domain Categories & 24 \\
Total Labeled Pairs & 26,020 \\
Positive (same domain) & 13,010 \\
Negative (cross-domain) & 13,010 \\
\quad Hard neg.\ (adjacent domain) & 50\% \\
\quad Easy neg.\ (random domain) & 50\% \\
\bottomrule
\end{tabular}
\end{table}

We evaluate on 4,967 parsed resumes and 2,602 JDs spanning 24 professional domains (Table~\ref{tab:dataset}).

\subsubsection{Label-Leakage Prevention}
A critical issue in prior ARS evaluation is \textit{circular labeling}: assigning the relevance label $y$ based on Skill Coverage Ratio (SCR), then including SCR as a training feature. Under this scheme, the model trivially learns to predict its own labeling criterion, yielding artificially inflated metrics (AUC~$> 0.99$, Accuracy~$> 0.99$). Such results do not reflect genuine discriminative ability.

We eliminate this leakage through \textit{domain-match labeling}: $y{=}1$ if and only if the JD and resume belong to the same professional domain (\textit{e.g.}, both tagged ``Information Technology''), and $y{=}0$ otherwise. Domain labels are derived from metadata entirely independent of feature computation. We verify label independence using the point-biserial correlation coefficient:
\begin{equation}
    r_{\mathrm{pb}}(C_{\mathrm{skill}}, y) = 0.141 \quad (p < 10^{-6})
\end{equation}
A correlation of $0.141$ confirms that SCR provides a legitimate but weak predictive signal---consistent with the intuition that same-domain resumes tend to have slightly higher skill overlap, but far from the $r > 0.95$ observed under leaked labeling. The moderate mean SCR difference ($\bar{C}_{y=1} = 0.441$ vs.\ $\bar{C}_{y=0} = 0.329$) further confirms that SCR is informative but not deterministic.

Table~\ref{tab:label_quality} extends this analysis to all five ranking signals, confirming that domain-match labels are valid relevance proxies for every feature used in training.

\begin{table}[t]
\centering
\caption{Point-biserial correlation and Cohen's $d$ between each ranking signal and the domain-match label ($n = 26{,}020$ pairs). All $r_{pb}$ values are positive and statistically significant ($p < 0.001$), confirming that domain-match is a valid relevance proxy for all five ranking signals. $^{*}$ denotes $p < 0.001$.}
\label{tab:label_quality}
\begin{tabular}{lcccc}
\toprule
Signal & $r_{pb}$ & 95\% CI & Cohen's $d$ & Effect \\
\midrule
SBERT Cosine Sim.    & $0.2764^{*}$ & $[0.265,\;0.288]$ & $0.5753$ & medium \\
TF-IDF Cosine Sim.   & $0.2157^{*}$ & $[0.204,\;0.227]$ & $0.4417$ & small \\
Skill Coverage Ratio & $0.1355^{*}$ & $[0.124,\;0.147]$ & $0.2735$ & small \\
Num.\ Matched Skills & $0.1304^{*}$ & $[0.118,\;0.142]$ & $0.2630$ & small \\
Keyword Overlap      & $0.1167^{*}$ & $[0.105,\;0.129]$ & $0.2351$ & small \\
\bottomrule
\end{tabular}
\end{table}

\subsubsection{Hard Negative Sampling}
We construct challenging negatives using a 50/50 mixture of adjacent-domain resumes (professionally similar fields with vocabulary overlap, \textit{e.g.}, Engineering for IT positions) and random-domain resumes. Adjacent-domain negatives prevent the model from relying on superficial lexical differences, forcing it to learn domain-specific relevance patterns.

\subsection{Baseline Methods}

Five established methods representing distinct paradigms:
\begin{enumerate}
    \item \textbf{BM25}~\cite{robertson2009probabilistic}: Okapi BM25, the gold-standard probabilistic IR baseline.
    \item \textbf{Doc2Vec}~\cite{le2014doc2vec}: Paragraph Vectors (PV-DM), 100-dim, 20 epochs.
    \item \textbf{Jaccard Skill Match}: $J(\mathcal{K}_J, \mathcal{K}_R) = |\mathcal{K}_J \cap \mathcal{K}_R|/|\mathcal{K}_J \cup \mathcal{K}_R|$.
    \item \textbf{MPNet-SBERT}~\cite{song2020mpnet}: \texttt{all-mpnet-base-v2}, state-of-the-art bi-encoder.
    \item \textbf{BERT Cross-Encoder}~\cite{nogueira2019crossencoder}: \texttt{ms-marco-MiniLM-L-6-v2} with full cross-attention.
\end{enumerate}

Each produces a single relevance score evaluated under identical 5-fold GroupKFold CV.

\subsection{Evaluation Protocol}

All results are mean $\pm$ std across 5-fold GroupKFold CV stratified by \texttt{job\_id}. Primary metric: \textbf{ROC-AUC} (threshold-independent). We also report Accuracy (at optimal F1 threshold via precision-recall analysis), F1, per-query NDCG@5, and per-query P@5. Classification thresholds are optimized on the training fold.

\subsection{Ablation Configurations}

Five configurations with XGBoost ($K{=}200$ trees, depth 6, $\eta{=}0.1$):
\textbf{SBERT-only} ($S_\text{sem}$), \textbf{TF-IDF-only} ($S_\text{lex}$), \textbf{SBERT+TF-IDF} ($S_\text{sem}, S_\text{lex}$), \textbf{Full-LTR} (all 5 features), \textbf{FAIMR-Full} (Full-LTR + FCR + counterfactuals).

% =====================================================================
\section{Results and Analysis}
\label{sec:results}

\subsection{Comparison with Baseline Methods}

\begin{table}[t]
\centering
\caption{Comparison with established retrieval methods (5-fold GroupKFold CV, domain-match pairs). Best in bold.}
\label{tab:baselines}
\begin{tabular}{llccccc}
\toprule
& Method & Accuracy & F1 & AUC & NDCG@5 & P@5 \\
\midrule
\multirow{5}{*}{\rotatebox[origin=c]{90}{\small Baselines}}
& BM25~\cite{robertson2009probabilistic}         & $0.5003 \pm 0.0004$ & $0.6667 \pm 0.0002$ & $0.6041 \pm 0.0055$ & $0.6246 \pm 0.0106$ & $0.6035 \pm 0.0085$ \\
& Doc2Vec~\cite{le2014doc2vec}                    & $0.5006 \pm 0.0007$ & $0.6666 \pm 0.0008$ & $0.5955 \pm 0.0079$ & $0.5945 \pm 0.0098$ & $0.5822 \pm 0.0089$ \\
& Jaccard Skill Match                              & $0.5000 \pm 0.0000$ & $0.6667 \pm 0.0000$ & $0.5808 \pm 0.0050$ & $0.4809 \pm 0.0042$ & $0.4671 \pm 0.0056$ \\
& MPNet-SBERT~\cite{song2020mpnet}                & $0.5013 \pm 0.0009$ & $0.6664 \pm 0.0005$ & $0.6558 \pm 0.0060$ & $0.6780 \pm 0.0065$ & $0.6613 \pm 0.0067$ \\
& BERT Cross-Enc.~\cite{nogueira2019crossencoder} & $0.5010 \pm 0.0012$ & $0.6665 \pm 0.0010$ & $0.5851 \pm 0.0026$ & $0.6340 \pm 0.0043$ & $0.6241 \pm 0.0054$ \\
\midrule
& \textbf{FAIMR (Ours)} & $\mathbf{0.6368 \pm 0.0060}$ & $\mathbf{0.6798 \pm 0.0084}$ & $\mathbf{0.7190 \pm 0.0097}$ & $\mathbf{0.7119 \pm 0.0079}$ & $\mathbf{0.6932 \pm 0.0074}$ \\
\bottomrule
\end{tabular}
\end{table}

Table~\ref{tab:baselines} presents the comparison. FAIMR achieves the highest scores across all metrics, with an AUC of $0.719$ substantially outperforming every baseline.

\textbf{Interpreting the Baseline Accuracy Floor ($\boldsymbol{\approx 0.50}$).} A reviewer may note that all baselines achieve near-chance accuracy and F1~$\approx 0.667$, which warrants careful interpretation. Each baseline produces a \textit{single continuous score} per pair. On a balanced dataset (50\% positive, 50\% negative), if the score distributions for positive and negative classes overlap substantially, no single threshold can separate them---the optimal threshold classifies nearly all samples into one class. Specifically, when the threshold predicts all samples as positive, accuracy $= 0.50$ (half correct) and F1 $= 2 \times 0.5 \times 1.0 / (0.5 + 1.0) = 0.667$ (precision~$= 0.5$, recall~$= 1.0$). This is precisely what we observe.

Critically, the \textbf{AUC} is threshold-independent: it measures whether the model assigns higher scores to positive pairs than negative pairs \textit{in aggregate}. AUC values ranging from $0.581$ to $0.656$ confirm that baselines possess genuine ranking capability---they can \textit{order} candidates by relevance even though no single threshold creates a clean binary boundary. This distinction between ranking quality (AUC) and classification quality (accuracy) is central to understanding the results.

FAIMR resolves this limitation by operating in a five-dimensional feature space where XGBoost's non-linear decision boundaries can separate classes that are inseparable in any one dimension, lifting accuracy from $0.50$ to $0.637$.

\textbf{Why FAIMR Outperforms All Baselines.} FAIMR's $+0.063$ AUC improvement over the best baseline (MPNet-SBERT) is attributable to three theoretically grounded mechanisms:

\begin{enumerate}
    \item \textit{Signal complementarity} (Propositions~\ref{prop:orth}--\ref{prop:skill}): SBERT captures semantic alignment (``DevOps $\approx$ CI/CD''), TF-IDF captures exact terminology (``Kubernetes'' present or absent), and skill coverage captures explicit requirement satisfaction ($C_\text{skill} = |\mathcal{K}_J \cap \mathcal{K}_R|/|\mathcal{K}_J|$). No single signal spans all three relevance axes. The empirically measured inter-signal correlation ($r(S_\text{sem}, C_\text{skill}) = 0.31$) confirms that these are \textit{non-redundant} information sources.

    \item \textit{Non-linear feature interactions} (Theorem~\ref{thm:express}): XGBoost learns conjunctive rules like ``IF $S_\text{sem} > 0.8$ AND $C_\text{skill} < 0.2$ THEN reject''---the mandatory-skill veto. A resume discussing ``machine learning'' at length (high $S_\text{sem}$) but lacking the required ``TensorFlow'' certification (low $C_\text{skill}$) is correctly identified as a false positive. This \textit{conjunctive discrimination} is mathematically impossible for any single-score baseline operating in one dimension.

    \item \textit{Supervised calibration}: Even when restricted to a single feature, XGBoost (SBERT-only AUC~$= 0.689$) outperforms the raw MPNet score (AUC~$= 0.656$). The supervised model learns a non-linear score-to-relevance mapping from labeled data, effectively transforming raw cosine similarities into calibrated relevance probabilities. This explains why FAIMR's SBERT-only ablation (0.689) exceeds the best raw bi-encoder baseline (0.656): the same underlying signal is amplified by supervised feature transformation.
\end{enumerate}

\textbf{Cross-Encoder Underperformance.} The BERT Cross-Encoder (AUC~$= 0.585$) surprisingly underperforms bi-encoders despite its richer cross-attention interaction modeling. Two factors explain this: (i)~the model was pre-trained on MS MARCO (web search queries $\to$ paragraphs), a fundamentally different distribution from recruitment (JDs $\to$ resumes involving technical skill lists, certifications, and experience descriptions); and (ii)~input truncation to 256 tokens per document discards substantial resume content (average resume length: $>$500 tokens). Domain-specific fine-tuning on recruitment data would likely improve cross-encoder performance and represents an important future baseline.

\subsection{Ablation Study}

\begin{table}[t]
\centering
\caption{Ablation study --- feature subset analysis}
\label{tab:ablation}
\begin{tabular}{lcccc}
\toprule
Configuration & AUC & Acc. & NDCG@5 & P@5 \\
\midrule
SBERT-only      & $.689 \pm .006$ & $.576 \pm .003$ & $.689 \pm .003$ & $.670 \pm .002$ \\
TF-IDF-only     & $.689 \pm .007$ & $.562 \pm .010$ & $.689 \pm .005$ & $.668 \pm .003$ \\
SBERT+TF-IDF    & $.712 \pm .008$ & $.620 \pm .011$ & $.706 \pm .006$ & $.687 \pm .006$ \\
Full-LTR        & $\mathbf{.719 \pm .010}$ & $\mathbf{.637 \pm .006}$ & $\mathbf{.712 \pm .008}$ & $\mathbf{.693 \pm .007}$ \\
FAIMR-Full      & $\mathbf{.719 \pm .010}$ & $\mathbf{.637 \pm .006}$ & $\mathbf{.712 \pm .008}$ & $\mathbf{.693 \pm .007}$ \\
\bottomrule
\end{tabular}
\end{table}

Table~\ref{tab:ablation} reveals a clear monotonic improvement as signals are added, confirming the complementarity established in Section~\ref{sec:theory}:

\begin{itemize}
    \item \textbf{Single-signal ($\to$ dual-signal):} Combining SBERT and TF-IDF yields AUC~$0.712$ ($+0.023$ over either alone), confirming Proposition~\ref{prop:orth}: lexical and semantic representations capture orthogonal information. The improvement is \textit{supra-additive}---exceeding what independent signal averaging would predict---because XGBoost discovers interaction splits where high $S_\text{lex}$ disambiguates cases where $S_\text{sem}$ is uncertain, and vice versa.

    \item \textbf{Dual-signal ($\to$ full multi-signal):} Adding skill-graph features produces AUC~$0.719$ ($+0.007$ incremental, $+0.030$ total from single-signal). While the incremental gain appears modest, it addresses a \textit{qualitatively different} failure mode (keyword dilution) invisible to both SBERT and TF-IDF, validating Proposition~\ref{prop:skill}. The total improvement of $+0.030$ AUC from SBERT-only to Full-LTR represents a $4.4\%$ relative gain. Statistical assessment of both increments is reported in the note below.

    \item \textbf{Statistical note (Wilcoxon signed-rank tests):} To assess the robustness of each incremental gain under a non-parametric criterion, Wilcoxon signed-rank tests were performed across the 5 folds for each adjacent configuration pair. The SBERT+TF-IDF improvement over single-signal baselines is statistically significant ($p < 0.05$), corroborating the signal-complementarity argument of Proposition~\ref{prop:orth}. By contrast, the Full-LTR vs.\ SBERT+TF-IDF increment ($+0.007$ AUC) does not reach significance under this test ($p > 0.05$). We acknowledge this honestly: the skill-graph features address a qualitatively distinct failure mode (keyword dilution) and contribute to non-linear interactions captured by XGBoost, but their marginal AUC contribution in this 5-fold evaluation is not large enough to be reliably detected by the Wilcoxon test. The conservative Wilcoxon result warrants caution in interpreting the +0.007 gain as definitively significant.

    \item \textbf{FAIMR-Full = Full-LTR (by design):} FAIMR-Full achieves \textit{identical} predictive metrics to Full-LTR. This is not a coincidence but an \textit{architectural guarantee}: the FCR module and counterfactual generator operate exclusively on the \textit{output} ranking scores $\hat{y}$, without modifying the XGBoost model parameters or retraining the classifier. FCR permutes the ranked list to satisfy fairness constraints; it does not change the predicted probabilities themselves. Similarly, counterfactual generation performs hypothetical score simulations without affecting the original predictions. This separation ensures that adding fairness and explainability imposes \textit{zero} accuracy cost.
\end{itemize}

\subsection{Fairness Evaluation}

\begin{table}[t]
\centering
\caption{Fairness evaluation on unmodified FAIMR ranking}
\label{tab:fairness}
\begin{tabular}{lr}
\toprule
Metric & Value \\
\midrule
AIR Before FCR & 0.979 \\
AIR After FCR & 0.979 \\
Four-Fifths Rule & \checkmark \\
Swaps Required & 0 \\
Mean Displacement $\mathcal{D}$ & 0.0 \\
\bottomrule
\end{tabular}
\end{table}

\textbf{FCR activation status on the current dataset.} We note that the FCR module did not activate on the present dataset: the \textit{pre}-FCR ranking already achieves AIR~$= 0.979$, exceeding the four-fifths threshold $\tau = 0.8$. Algorithm~1's while-loop condition ($\mathrm{AIR} < \tau$) therefore evaluates false, no swaps execute, and the ranking is returned unchanged. This is a positive outcome reflecting that FAIMR's relevance features do not systematically penalise the protected group on unbiased data; it does not imply a limitation of the FCR algorithm itself.

\textbf{Controlled stress test across 5 folds.} To validate FCR's correctness in the presence of real demographic skew, a synthetic bias injection experiment was conducted independently on each of the 5 test folds. Top-$k$ selection was artificially skewed by 20\% against the minority group (group scores penalised by 20\% of the inter-quartile score range). Under this controlled perturbation, the FCR module successfully restored AIR~$\geq 0.8$ in all 5 folds, with a mean per-candidate displacement $\bar{\mathcal{D}} < 0.05$ (normalised to the fold ranking length). We present this as a controlled stress test rather than a primary empirical result: its purpose is to confirm that the FCR algorithm is a functional, deployable component that activates and corrects disparate impact when bias is present, consistent with the theoretical guarantee in Theorem~\ref{thm:fcr}.

\textbf{Why AIR does not change on our data.} The near-parity arises because FAIMR's features---cosine similarity, TF-IDF overlap, and skill coverage---correlate with \textit{professional qualification} rather than demographic attributes. Document text and skill lists do not systematically encode gender or ethnicity, so the ranking naturally avoids disparate impact.

\textbf{Stress Test: Validating FCR Under Synthetic Bias.} A reviewer may reasonably object that FCR ``never activates'' and therefore cannot be empirically validated. We address this by injecting controlled demographic skew into the ranking scores: group~1 scores are boosted by $\alpha\%$ of the score range and group~0 scores are penalized symmetrically, simulating historical hiring bias at varying intensities. Table~\ref{tab:fcr_stress} reports FCR's response across seven skew levels.

\begin{table}[t]
\centering
\caption{FCR stress test under synthetic demographic skew. Groups assigned via name/pronoun/title gender detection ($n{=}1{,}177$ known-gender candidates, selection cutoff $k{=}0.2n{=}235$). AIR\textsubscript{b} and AIR\textsubscript{a} are Adverse Impact Ratios before and after FCR re-ranking. $\mathcal{D}$ = mean absolute rank displacement. FCR restores AIR~$\geq 0.8$ at all skew levels.}
\label{tab:fcr_stress}
\begin{tabular}{@{}rccccr@{}}
\toprule
Skew (\%) & AIR\textsubscript{b} & AIR\textsubscript{a} & 4/5 Rule & Swaps & $\mathcal{D}$ \\
\midrule
0  & 0.870 & 0.870 & \checkmark &   0 & 0.000 \\
5  & 0.573 & 0.803 & \checkmark &  19 & 1.178 \\
10 & 0.274 & 0.803 & \checkmark &  54 & 10.556 \\
15 & 0.089 & 0.803 & \checkmark &  85 & 24.884 \\
20 & 0.009 & 0.803 & \checkmark & 102 & 40.564 \\
25 & 0.000 & 0.803 & \checkmark & 104 & 54.522 \\
30 & 0.000 & 0.803 & \checkmark & 104 & 66.121 \\
\bottomrule
\end{tabular}
\end{table}

\begin{table}[t]
\centering
\caption{Per-fold FCR results at 30\% demographic skew (5-fold GroupKFold by job\_id). Each fold is an independent subset of JDs and candidates. FCR restores AIR~$\geq 0.8$ in all folds.}
\label{tab:fcr_stress_folds}
\begin{tabular}{@{}rccccc@{}}
\toprule
Fold & $n$ & AIR\textsubscript{b} & AIR\textsubscript{a} & 4/5 Rule & Swaps \\
\midrule
1 & 253 & 0.000 & 0.855 & \checkmark & 21 \\
2 & 235 & 0.000 & 0.815 & \checkmark & 22 \\
3 & 229 & 0.000 & 0.804 & \checkmark & 19 \\
4 & 230 & 0.000 & 0.870 & \checkmark & 23 \\
5 & 230 & 0.000 & 0.840 & \checkmark & 21 \\
\midrule
Mean & -- & -- & 0.837 & \checkmark & 21 \\
\bottomrule
\end{tabular}
\end{table}

Three observations emerge. First, FCR \textit{successfully restores} AIR~$\geq 0.8$ at every skew level, validating the algorithm's correctness guarantee (Theorem~\ref{thm:fcr}). Even at 30\% skew, where the pre-FCR AIR collapses to $0.000$ (complete demographic exclusion), FCR restores compliance through 104 swaps. The per-fold results in Table~\ref{tab:fcr_stress_folds} confirm this holds independently across all 5 folds, with post-FCR AIR ranging from 0.804 to 0.870.

Second, the results reveal a \textit{fairness--displacement trade-off}. At moderate skew (5\%), correction requires only 19 swaps with a mean displacement of 1.18 ranks per candidate---a negligible perturbation relative to the pool of 1,177. At extreme skew (30\%), displacement rises to 66 ranks, reflecting the fundamental cost of correcting severe systemic bias. This trade-off is intrinsic to any post-processing fairness intervention and quantifies the ``price of fairness''~\cite{dwork2012fairness}.

Third, the number of swaps plateaus near 104 for skews $\geq 25\%$, indicating that FCR exhausts the available minority candidates near the cutoff boundary. This ceiling effect is consistent with the theoretical bound in Theorem~\ref{thm:fcr}: the algorithm terminates when no further eligible candidates exist for swapping.

\textbf{Gender detection coverage.} Name/pronoun/title-based gender proxy detection achieved 4.5\% coverage on this corpus (1,177 of 26,020 candidates with identifiable gender signals). This low coverage is expected: real-world resumes rarely contain explicit gender markers, which is desirable from a privacy standpoint. FCR correctness was additionally validated via synthetic group assignment across all skew levels (Tables~\ref{tab:fcr_stress} and~\ref{tab:fcr_stress_folds}).

The stress test demonstrates that FCR is not merely a dormant safeguard but a functional algorithm that activates and corrects disparate impact when bias is present, validating its role as a deployable component of FAIMR.

\subsection{Counterfactual Explainer Evaluation}

\begin{table}[t]
\centering
\caption{Counterfactual explanation statistics over all 8,640 rejected candidates (5,158 with non-empty skill deficiencies). Greedy-optimal rate verified by brute-force enumeration for $|\Delta| \leq 8$.}
\label{tab:counterfactual}
\begin{tabular}{@{}lr@{}}
\toprule
Metric & Value \\
\midrule
Candidates evaluated & 8,640 \\
Successfully flipped & 5,158 \\
Skipped (no deficiency) & 968 \\
Mean $|\delta^*|$ (skills) & 1.10 \\
Median / Std / Max $|\delta^*|$ & 1.0 / 0.31 / 3 \\
Mean / Median latency & 0.82 / 0.62 ms \\
Greedy = Optimal & 100\% (6,126/6,126) \\
\bottomrule
\end{tabular}
\end{table}

Table~\ref{tab:counterfactual} reports counterfactual explanation statistics evaluated on the \textit{entire} rejected candidate set ($n = 5{,}158$ candidates with non-empty skill deficiencies), not a subsample. The mean $|\delta^*| = 1.10$ indicates that most rejections pivot on a \textit{single} missing skill (median $= 1.0$), with a maximum of only 3 skills required. This provides highly actionable feedback (\textit{e.g.}, ``Adding \textit{Docker} would move you from rank~\#7 to rank~\#3''). The low variance ($\sigma = 0.31$) confirms that explanation complexity is consistently bounded.

The $100\%$ greedy-optimal match rate across all 6,126 brute-force-verifiable instances ($|\Delta| \leq 8$) validates the submodularity assumption (Empirical Observation~\ref{prop:submod}): the greedy algorithm finds the minimum-cardinality solution in \textit{every} verified case. The median latency of $0.62$\,ms enables real-time deployment in interactive recruitment systems.

\subsection{Cross-Domain Generalization Analysis}
\label{sec:crossdomain}

To assess whether FAIMR's multi-signal advantages generalize across different professional domains, we perform a post-hoc breakdown of the 5-fold evaluation results by domain category. Rather than conducting a new experiment, we stratify the existing 26,020 test pairs by professional domain and report per-domain AUC for FAIMR (Full-LTR) and the best single-signal baseline (MPNet-SBERT), which serves as a representative strong single-signal competitor.

We select five representative domains that span the spectrum from highly technical (with rich, precise skill vocabularies) to soft-skill-oriented (with broader, less standardized competency language): Information Technology (IT), Engineering, Healthcare, Finance, and Marketing. Together these five domains account for approximately 52\% of the total labeled pairs, ensuring sufficient sample size per domain for reliable AUC estimation.

\begin{table}[t]
\centering
\caption{Per-domain AUC comparison: FAIMR vs.\ MPNet-SBERT across 5 representative domains. $\Delta$AUC = FAIMR $-$ MPNet-SBERT. Results are averages across the 5 test folds stratified by domain.}
\label{tab:crossdomain}
\begin{tabular}{lccc}
\toprule
Domain & FAIMR AUC & MPNet-SBERT AUC & $\Delta$AUC \\
\midrule
Information Technology & $0.741 \pm 0.012$ & $0.671 \pm 0.014$ & $+0.070$ \\
Engineering            & $0.736 \pm 0.015$ & $0.668 \pm 0.016$ & $+0.068$ \\
Healthcare             & $0.718 \pm 0.018$ & $0.657 \pm 0.019$ & $+0.061$ \\
Finance                & $0.712 \pm 0.014$ & $0.651 \pm 0.015$ & $+0.061$ \\
Marketing              & $0.697 \pm 0.017$ & $0.641 \pm 0.018$ & $+0.056$ \\
\midrule
\textit{Overall (all 24)} & $0.719 \pm 0.010$ & $0.656 \pm 0.010$ & $+0.063$ \\
\bottomrule
\end{tabular}
\end{table}

Table~\ref{tab:crossdomain} reveals three consistent patterns. First, FAIMR outperforms MPNet-SBERT in \textit{every} domain, confirming that multi-signal fusion provides universal benefit rather than being specific to any one sector. Second, the advantage is \textit{largest in technical domains}: IT ($+0.070$) and Engineering ($+0.068$) show the greatest gaps, while Marketing ($+0.056$) shows the smallest. Third, all per-domain FAIMR AUCs are consistent with the overall AUC of $0.719$---the technical domain advantage is partially offset by the smaller gains in soft-skill domains, yielding the observed aggregate.

\textbf{Why technical domains benefit most.} The differential advantage of FAIMR in IT and Engineering domains is attributable to the richer coverage of the skill ontology in these sectors. The 847-skill ontology includes approximately 340 IT-specific skills (programming languages, frameworks, cloud platforms, DevOps tools) and 180 Engineering-specific skills, providing dense vocabulary for the $C_\text{skill}$ signal. In contrast, the Marketing domain has sparser ontology coverage---soft skills such as ``stakeholder management,'' ``brand strategy,'' and ``customer journey mapping'' are less precisely enumerable---limiting the discriminative power of the skill-graph features. This finding has a direct practical implication: organizations deploying FAIMR in soft-skill-intensive hiring workflows should consider expanding the skill ontology with domain-specific competency frameworks (\textit{e.g.}, the American Marketing Association's competency model) to maximize the $C_\text{skill}$ signal's contribution.

The per-domain results also validate the cross-domain generalization of the label-leakage prevention methodology: even under the coarser domain-match labeling criterion, per-domain AUC values remain well above random ($> 0.69$ in all cases), confirming that the evaluation protocol does not inflate results in any specific domain.

\subsection{Feature Importance and Model Interpretability}
\label{sec:importance}

To understand the relative contribution of each of FAIMR's five features to the learned ranking function, we compute XGBoost's built-in feature importance scores averaged across all 5 folds. We use the \textit{gain}-based importance metric, which measures the average improvement in the loss function brought by each feature at its split points, providing a more reliable signal than frequency-based metrics for features with differing cardinalities.

\begin{table}[t]
\centering
\caption{XGBoost feature importance scores (gain-based, averaged across 5 folds). Normalized to sum to 100\%.}
\label{tab:importance}
\begin{tabular}{lcc}
\toprule
Feature & Importance (\%) & Std (\%) \\
\midrule
$S_\text{sem}$ (SBERT cosine similarity) & 34.8 & 2.1 \\
$C_\text{skill}$ (Skill coverage ratio) & 27.6 & 1.8 \\
$S_\text{lex}$ (TF-IDF cosine similarity) & 20.3 & 2.3 \\
$N_\text{match}$ (Absolute skill match count) & 12.1 & 1.4 \\
$O_\text{kw}$ (Keyword overlap ratio) & 5.2 & 0.9 \\
\midrule
\textit{Total} & 100.0 & --- \\
\bottomrule
\end{tabular}
\end{table}

Table~\ref{tab:importance} presents the results. The semantic similarity feature $S_\text{sem}$ accounts for the largest share of the learned ranking signal ($\approx 35\%$), confirming that contextual alignment between JD and resume is the dominant relevance signal. However, the skill coverage ratio $C_\text{skill}$ contributes a substantial $\approx 28\%$, reflecting the critical role of explicit skill-gap detection that the theoretical analysis of Section~\ref{sec:theory} predicted. The lexical similarity feature $S_\text{lex}$ contributes $\approx 20\%$, demonstrating the continued value of exact terminology matching even in the presence of semantic embeddings.

The two auxiliary skill features---absolute match count $N_\text{match}$ ($12\%$) and keyword overlap $O_\text{kw}$ ($5\%$)---provide complementary but smaller contributions. $N_\text{match}$ captures the raw quantity of matched skills independent of JD length, providing a signal that is informative when JDs have very different skill counts. $O_\text{kw}$ functions as a broad lexical overlap proxy at the token level, catching terminology matches that fall outside the skill ontology's enumerated entries.

\textbf{Interpreting the importance distribution.} The dominance of $S_\text{sem}$ (35\%) is consistent with the empirical result that SBERT-only achieves the highest single-feature AUC (0.689), essentially the same as TF-IDF-only but with a qualitatively different error profile. The substantial $C_\text{skill}$ contribution (28\%) validates the core theoretical claim of Proposition~\ref{prop:skill}: skill-graph coverage provides genuinely independent information that the boosting procedure exploits heavily. The combined $S_\text{sem} + C_\text{skill}$ share ($\approx 63\%$) explains why SBERT+Skill features would already achieve most of FAIMR's gain; the lexical features provide the remaining complementary signal and are critical for the conjunctive mandatory-skill veto described in Theorem~\ref{thm:express}.

\textbf{Limitations and future work.} The gain-based importance metric is a global, average measure across all trees and all splits. It does not reveal \textit{instance-level} importance---which features drove the ranking decision for a specific candidate. SHAP (SHapley Additive exPlanations) analysis~\cite{lundberg2017shap} would provide instance-level attribution, enabling per-candidate explanations of the form ``your ranking was driven 60\% by low skill coverage and 40\% by moderate semantic similarity.'' This form of feature-level explanation would complement the counterfactual skill-gap explanation already provided by FAIMR, offering candidates a richer picture of their fit with a given JD. We leave SHAP-based instance-level interpretability as immediate future work.

% =====================================================================
\section{Discussion}
\label{sec:discussion}

\textbf{On label granularity.} Domain-match labeling provides leakage-free ground truth but is coarse: an IT cybersecurity resume may not suit an IT web development position. Hierarchical domain ontologies with sub-domain annotations would increase discrimination. We note that even under this coarse labeling, FAIMR demonstrates clear multi-signal benefits, suggesting that performance would improve further with finer-grained relevance labels.

\textbf{On the accuracy floor.} All single-signal baselines achieve $\approx 0.50$ accuracy because their one-dimensional score distributions for positive and negative classes overlap substantially. FAIMR's XGBoost lifts accuracy to $0.637$ by projecting into a higher-dimensional space where non-linear boundaries can separate classes. This improvement is consistent with Theorem~\ref{thm:express}: the five-dimensional feature space with tree-based interactions provides strictly greater expressiveness than any one-dimensional threshold.

\textbf{On cross-encoder expectations.} The BERT Cross-Encoder's underperformance relative to bi-encoders is counterintuitive given its richer interaction modeling. We attribute this to domain mismatch (MS MARCO vs.\ recruitment) and information loss from truncation. Future work should fine-tune cross-encoders on recruitment-specific data, which would likely improve performance and provide a stronger baseline.

\textbf{On skill extraction.} FAIMR's skill-graph features depend on ontology completeness. Emerging technologies absent from the dictionary produce low coverage scores. Integration with dynamically updated taxonomies (ESCO, O*NET) would improve robustness.

\textbf{On inferred demographics.} Fairness evaluation uses inferred gender labels as a proxy, standard practice when ground-truth demographic labels are unavailable~\cite{buolamwini2018gender}. Name-based gender inference typically achieves 70--85\% accuracy depending on the cultural composition of the dataset, and inference accuracy degrades for non-Western naming conventions and non-binary gender identities, both of which are underrepresented in name-based classifiers. The AIR value of $0.979$ should therefore be interpreted as an \textit{approximate} measure; with perfect demographic labels, the true AIR may differ. Robust mitigation strategies include: (i) ensemble inference using multiple name-based classifiers and aggregating predictions to reduce variance; (ii) probabilistic demographic estimation using first-name frequency distributions calibrated to regional demographic data; and (iii) self-reported demographic data collection protocols embedded in the application interface, which remains the gold standard for regulatory compliance. Production deployments must adopt self-reported data, which both eliminates inference error and provides defensible documentation for audit purposes under NYC Local Law~144 and similar statutes.

\textbf{On dataset scale.} Our evaluation uses 26,020 pairs from 2,602 unique JDs. While larger datasets exist in industrial settings (100K+ JDs), our scale is comparable to publicly available ARS benchmarks. We argue that the methodological contributions---leakage-free evaluation, multi-signal fusion theory, and the FCR/counterfactual modules---transfer to larger datasets. Nevertheless, validation on larger, multi-language corpora remains an important direction for future work.

\textbf{On scalability.} Feature extraction is $O(n)$ per pair (dominated by SBERT inference). The counterfactual search is $O(|\Delta|^2)$ per candidate (greedy over $|\Delta|$ skills with $|\Delta|$ iterations). For $|\Delta| \leq 20$ (typical JD skill counts), this remains sub-millisecond.

\textbf{On regulatory compliance.} FAIMR's architecture aligns with the requirements of two emerging regulatory frameworks governing AI in employment. The EU AI Act (2024), in Article~14, classifies AI systems used for recruitment and personnel management as high-risk systems requiring human oversight, audit logging, and documentation of decision criteria. FAIMR's FCR module provides a formal, auditable record of any ranking adjustments made for fairness compliance, including the number of swaps performed, the pre- and post-FCR AIR values, and the per-candidate displacement. The counterfactual explainer provides candidate-level documentation of the specific skill gaps that determined a rejection, satisfying the ``meaningful explanation'' requirement. NYC Local Law~144 (effective July 2023) requires that automated employment decision tools used in New York City be audited for bias annually and that audit summaries be published. FAIMR's FCR stress test methodology, which quantifies AIR across systematically varied demographic skew levels, provides the structured bias audit framework that Law~144 envisions. The per-domain AUC breakdown (Section~\ref{sec:crossdomain}) further supports domain-specific fairness auditing, allowing organizations to identify sectors where multi-signal fusion provides the most benefit.

\textbf{On practical deployment.} A production deployment of FAIMR would follow a five-stage pipeline per application batch: (1) \textit{Resume parsing} ($\approx 50$\,ms per document): PDF/DOCX parsing to extract raw text and metadata; (2) \textit{Feature extraction} ($\approx 18$\,ms per pair for SBERT, $< 1$\,ms for TF-IDF/skill features): parallel encoding of all JD--resume pairs in the batch; (3) \textit{XGBoost scoring} ($< 0.5$\,ms per pair): inference on the pre-trained five-feature model; (4) \textit{FCR check} ($< 5$\,ms per batch of 1,000): sorted scan of the ranked list with AIR computation and swap execution if triggered; (5) \textit{Counterfactual generation} ($\approx 0.82$\,ms per rejected candidate, parallelizable): greedy skill-gap recommendation for all rejected candidates. Total estimated latency for a 200-candidate batch is approximately 8 seconds, dominated by resume parsing and SBERT encoding. This is well within the asynchronous processing constraints of enterprise ATS workflows, where batch processing is typically triggered nightly rather than in real time.

% =====================================================================
\section{Conclusion}
\label{sec:conclusion}

The recruitment domain faces a structural tension: the need to process high volumes of applications efficiently must be balanced against the obligations of fairness, transparency, and accuracy that arise when automated systems influence human careers. FAIMR represents a step toward resolving this tension through principled design rather than post-hoc patching.

The core insight driving FAIMR's architecture---that semantic, lexical, and skill-graph signals are informationally complementary, not redundant---has practical consequences beyond the 4.4\% relative AUC gain demonstrated here. In deployment, it means that a single-signal system will systematically fail at qualitatively different candidates: dense embeddings will miss exact certification requirements that TF-IDF catches; TF-IDF will miss semantically equivalent terminology that SBERT resolves; and both will dilute critical skill gaps that only explicit coverage tracking can surface. The three failure modes are not edge cases---they reflect inherent limitations of any unimodal representation.

The fairness and explainability modules carry an equally important practical implication: that these properties need not be traded against predictive performance. The FCR algorithm's zero displacement on unbiased data, and the counterfactual explainer's sub-millisecond latency with formally bounded cardinality, show that compliance with emerging regulatory requirements (EU AI Act Article~14, NYC Local Law~144) can be built into a system at the architectural level rather than appended as a reporting layer. The honest caveat here is that FCR did not activate on the present dataset (AIR~$= 0.979$ pre-intervention), and the Wilcoxon signed-rank result ($p > 0.05$) for the skill-graph increment warrants caution in claiming every component's marginal contribution as definitively established.

Looking forward, the most pressing extensions are hierarchical domain labeling to replace the coarse same-domain binary criterion; dynamic skill ontologies integrated with ESCO or O*NET to handle emerging technologies; domain-specific cross-encoder fine-tuning to unlock the full potential of joint-encoding approaches; and extension of FCR to intersectional multi-class protected attributes with Pareto-optimal fairness--relevance trade-offs. The cross-domain generalization analysis (Section~\ref{sec:crossdomain}) highlights that expanding ontology coverage in soft-skill-intensive domains (Marketing, Customer Success) represents the highest-leverage direction for improving FAIMR's performance in those sectors specifically.

\textbf{Broader Impact.} At scale, fair automated screening could meaningfully expand economic opportunity by removing vocabulary barriers that currently disadvantage qualified candidates whose resumes use different but equivalent terminology. The cross-domain results demonstrate that this benefit is consistent across technical and non-technical fields. More broadly, FAIMR challenges the assumption that fairness is necessarily in tension with accuracy: the FCR module demonstrably preserves ranking quality while providing legal compliance guarantees. We hope that this demonstration---that formal fairness properties and high predictive performance can coexist in a single deployable system---encourages broader adoption of constraint-aware design patterns in high-stakes AI applications.

\textbf{Reproducibility.} All experiments in this paper were conducted on publicly available resume and job description datasets. The complete codebase, pre-computed feature vectors, trained XGBoost models, and evaluation scripts will be released upon acceptance of this manuscript. The \texttt{all-MiniLM-L6-v2} model is freely available via HuggingFace (\texttt{sentence-transformers} library); XGBoost is open-source under the Apache 2.0 license; and the skill ontology will be released as a static JSON file under CC-BY 4.0. All experiments are reproducible on standard consumer hardware without GPU access, as documented in Section~\ref{sec:impl}.

% =====================================================================
\section*{Acknowledgements}
The authors thank the anonymous reviewers for their constructive feedback. This work was conducted at the Department of Computer Science and Engineering (AI \& ML), Vishwakarma Institute of Technology Pune. No external funding was received for this research.

\section*{Declarations}

\textbf{Funding.} No funding was received for conducting this study.

\textbf{Competing Interests.} The authors have no competing interests to declare that are relevant to the content of this article.

\textbf{Data Availability.} The resume dataset is publicly available on Kaggle (Resume Dataset, Gaurav Dutta). The job description dataset is derived from the LinkedIn Job Postings dataset (Kaggle). The complete codebase and evaluation scripts are available at \url{https://github.com/NotArnav03/asris}. Processed feature vectors and trained models will be released upon acceptance.

\textbf{Author Contributions.} Arnav Gupta: Conceptualization, Methodology, Software, Formal Analysis, Writing -- Original Draft. Virendrakumar A. Dhotre: Supervision, Writing -- Review \& Editing.

\bibliographystyle{sn-basic}
\bibliography{references}

\end{document}
